\documentclass[output=paper]{langscibook}
\ChapterDOI{10.5281/zenodo.22880272}
\title{\textit{O aga} – The outrigger float}
\author{Ruth Sima Rigamu}
\abstract{}
\begin{document}
\maketitle\label{chap:23}
\begin{figure}[H]
    \captionsetup[subfigure]{oneside,margin={0cm,0cm}}
    \subfloat[]{\includegraphics[trim=0 0 0 0,clip,width=.34\linewidth]{figures/23_1.jpg}\label{fig:23_1}}\hfil
    \subfloat[]{\includegraphics[trim=0 0 0 0,clip,width=.3\linewidth]{figures/23_2.png}\label{fig:23_2}}\hfil
    \subfloat[]{\includegraphics[trim=0 0 0 0,clip,width=.35\linewidth]{figures/23_3.jpg}\label{fig:23_3}}
    \caption{(a) Young \textit{kopihoo}, (b) Adult \textit{kopihoo}, (c) \textit{Kopihoo} leaf}
\end{figure}
\begin{xltabular}{\textwidth}{lXX}
    (\ref{ex:text17-1})&\textit{Kopihoo}&\textit{Kopihoo} tree\\
    (\ref{ex:text17-2})&\textit{O naono vai o kopihoo na paku riori bono aga no sinivi.}&The \textit{kopihoo} tree is used for making the float of the outrigger canoe.\\
    (\ref{ex:text17-3})&\textit{A abana na mosi maari bono kopihoo vai o totobin, ore paa hihi maa bari rasuu, ore paa dee tavus vo hanana.}&The men cut a straight \textit{kopihoo} tree, then they remove its bark in the jungle and carry it out of the bush to a trail.\\
    (\ref{ex:text17-4})&\textit{Vaasun vai ori, ore paa raka.}&They stand it up there to dry.\\
    (\ref{ex:text17-5})&\textit{Be raka, ore paa dee voosu maa bari,}&When it is dry, they carry it home\\
    (\ref{ex:text17-6})&\textit{ore paa ma taraha bona ihu no kopihoo.}&and plane the surface of the nose of the \textit{kopihoo} log.\\
    (\ref{ex:text17-7})&\textit{Ore paa taraha bona meha papana.}&Then they plane the surface of the other side.\\
    (\ref{ex:text17-8})&\textit{Be kavara, ore paa asi ni bono buo ihu aga ore tara vamataa.}&After that, they singe the two noses of the float to look beautiful.\\
    (\ref{ex:text17-9})&\textit{Bea abana pampamu vakavara ore paa tapa ni bono aga teo sinivi.}&When the men have done the lashing (i.e. have lashed the \textit{vanai} ‘boom’ to the canoe), they fix the float to (the boom of) the canoe with pegs.\\
    (\ref{ex:text17-10})&\textit{O buo naono vai to tapa aga ni raara ei o kopihoo bara o roge.}&The two trees that we build floats from, these are the \textit{kopihoo} tree and the coastal hibiscus.\\
\end{xltabular}
\begin{figure}[H]
    \includegraphics[trim=0 0 0 0,clip,width=.6\linewidth]{figures/23_4.jpg}\hfil
    \includegraphics[trim=0 0 0 0,clip,width=.4\linewidth]{figures/23_5.jpg}\hfil
    \includegraphics[trim=0 0 0 0,clip,width=.6\linewidth]{figures/23_6.jpg}
    \caption{\textit{O aga no sinivi} -- The float of the indigenous canoe}
    \label{fig:23_4-6}
\end{figure}
\end{document}